\label{der:2-3}
\textbf{秩的三种定义一致。}可逆行变换保列间关系，因$EA\vx=0\Longleftrightarrow A\vx=0$；列变换把列组作可逆线性组合，保列空间维数。对行空间反向论证也成立。取任一非零元作主元，换到左上角，归一并用行、列倍加清掉同列同行的其余元，对右下块重复，得到$PAQ=\diag(I_r,0)$的矩形版本。该标准形行、列独立方向都恰为$r$。

还须说明非零子式的最大阶数未变。对任意一次行操作，任一$k$阶新子式按被改变的行多线性展开，是旧矩阵$k$阶子式的线性组合（重复行项为零）；故若全部旧$k$阶子式为零，新子式也全零。逆行操作给反向蕴含；列操作同理。因此最大非零子式阶数不变，标准形中恰为$r$。行秩、列秩、子式秩由此相等。

\textbf{乘积与Gram秩。}$\im AB\subseteq\im A$给$\rank AB\le\rank A$；映射$A$把$\im B$送到$\im AB$，输出维数不超过输入维数，给$\rank AB\le\rank B$。实数域中，$A\vx=0$直接推出$A\T A\vx=0$；反向若$A\T A\vx=0$，则
\[0=\vx\T A\T A\vx=\sum_i(A\vx)_i^2,\]
每项非负，所以每项为零。两矩阵核相同，未知量数相同，故秩相同。

\textbf{伴随秩三类。}若$\rank A=n$，$\adj A=(\det A)A\inv$可逆，秩$n$。若$\rank A=n-1$，存在非零$(n-1)$阶子式，所以$\adj A\ne0$；又$A\adj A=0$，伴随每列在一维$\ker A$中，故其秩至多1且非零，只能为1。若$\rank A\le n-2$，全部$(n-1)$阶子式为零，所以伴随为零。这里$n\ge2$确保余子式阶数至少1。
